\section{Valuative optimization and finite windows}
\label{sec:valuative}

We now prepare the proof of Theorem~\ref{thm:phase}. Throughout this and
the next two sections, put $k_0=\C$, $L=k_0(V)$, and $\eta=\theta_V$.
We use an invariant affine representative of the pointed klt Gorenstein
germ, on which $\eta$ is a semi-invariant canonical generator. The cyclic
group is denoted by $G$. All valuations in these sections are trivial on
$k_0$. Constants described as uniform are uniform along the subsequence
already chosen in Proposition~\ref{prop:product}.

Let $v$ be the invariant rational divisorial valuation of
Proposition~\ref{prop:cone-estimates}, normalized by $A_\eta(v)=1$.
Choose centred regular eigenfunctions $h_i$ as in that proposition,
including centred affine generators, and put $v_i=v(h_i)>0$. Thus
\begin{equation}\label{val:comparison}
 s(w):=\min_i\frac{w(h_i)}{v_i},\qquad
 w(g)\ge s(w)v(g),\qquad
 w(g)\le 2nA_\eta(w)v(g)
\end{equation}
for nonzero regular $g$ and centred quasi-monomial $w$. Positivity on
the centred affine generators ensures that the centre is $o$.
Choose centred regular eigenfunctions $f_1^0,\ldots,f_n^0$ whose
$v$-initials are algebraically independent and attain the successive
homogeneous transcendence ranks. Set $d_j=v(f_j^0)$, in nondecreasing
order. The preceding propositions give
\begin{equation}\label{val:degree-data}
 d_j\ge\delta>0,\qquad d_1\le C,\qquad
 \vol(v)\prod_{j=1}^n d_j\le C.
\end{equation}
The functions may be chosen as germs and then regarded as rational
functions; shrinking the affine representative accommodates any finite
list subsequently used.

\subsection{A bound on every finite window}

\begin{lemma}\label{val:window}
Put $d_{n+1}=+\infty$. There is a uniform constant $C_W$ such that
\begin{equation}\label{val:window-bound}
 \mathcal H(T):=\dim_{k_0}\cO_{V,o}/(v>T)
 \le C_W\prod_{j\le m}\frac{T}{d_j}
 \quad(d_m\le T<d_{m+1},\ 1\le m\le n).
\end{equation}
\end{lemma}

\begin{proof}
The restriction of $v$ to $k_0(f_1^0,\ldots,f_n^0)$ is the weighted
Gauss valuation. The weight-zero Laurent monomials for a basis of the
kernel of $\Z^n\to\Q$, $\alpha\mapsto\sum\alpha_jd_j$, have
algebraically independent residues. Refine $v$ by a full-rank monomial
valuation on these residues, and extend this refinement to the finite
extension of residue fields upstairs. We obtain a composite valuation
$\tau$ with residue field $k_0$. Using $\tau(f_j^0)$ as coordinates,
its value group is a lattice $\Gamma\subset\Q^n$ containing $\Z^n$
with finite index, and the original real value is the pairing with
$(d_1,\ldots,d_n)$.

Suppose centred regular functions $l_1,\ldots,l_n$ have independent
$\tau$-values generating a sublattice $\Lambda$. Every class of
$\Gamma/\Lambda$ has a representative which is the value of a regular
function: the regular value semigroup generates $\Gamma$, and its image
in a finite group is a subgroup. Fix one such representative per class.
Multiplying by the monomials $l^\nu$, $\nu\in\Z_{\ge0}^n$, gives
distinct $\tau$-values. Those of $v$-value at most $T$ are independent
modulo $(v>T)$. Comparing the leading coefficients as $T\to\infty$
gives
\begin{equation}\label{val:lattice-index}
 [\Gamma:\Lambda]\le \vol(v)\prod_{i=1}^n v(l_i).
\end{equation}
The fixed representatives contribute only bounded shifts in this
asymptotic comparison; no uniform bound for those shifts is needed.

Applying \eqref{val:lattice-index} to $f^0$ gives
$[\Gamma:\Z^n]\le C$. If the $j$th coordinate $\tau(g)_j$ of a
regular nonunit $g$ is nonzero, replace $f_j^0$ by $g$. The determinant
of the new lattice basis is $|\tau(g)_j|$, so
\[
 [\Gamma:\Z^n]|\tau(g)_j|
 \le \vol(v)\,v(g)\prod_{i\ne j}d_i.
\]
Consequently
\begin{equation}\label{val:lattice-box}
 |\tau(g)_j|\le C\frac{v(g)}{d_j}.
\end{equation}
The inequality is automatic when that coordinate is zero. Units have
$\tau$-value zero.

If $v(g)\le T<d_{m+1}$, the initial of $g$ is algebraic over the
graded fraction field generated by the initials of
$f_1^0,\ldots,f_m^0$. Take a homogeneous polynomial relation. On
refining its values by $\tau$, at least two lowest terms tie, and they
involve distinct powers of the initial of $g$: coefficients belong to
the graded field generated by an independent tuple, where their
nonzero initials cannot cancel. Thus $\tau(g)$ lies in the rational
span of the first $m$ coordinate vectors. The number of lattice points
in this subspace satisfying \eqref{val:lattice-box} and $v(g)\le T$
is at most
$C\prod_{j\le m}(1+T/d_j)\le C'\prod_{j\le m}(T/d_j)$.
Here one can place $\Gamma$ in
$\frac1e\Z^n$ with $e=[\Gamma:\Z^n]\le C$.
Finally $\tau$ has one-dimensional leaves. The bounded box has only
finitely many lattice points, so successively taking initials in the
finite jet quotient bounds its dimension by this count. This proves
\eqref{val:window-bound}.
\end{proof}

\subsection{Gauss valuations and continuation}

We say that a valuation is \emph{Gauss on a tuple} when the tuple has
algebraically independent homogeneous initials; equivalently, every
nonzero polynomial in the tuple has the minimum of its monomial
weights. We will also use this definition when the assigned weights
have rational relations.

\begin{lemma}\label{val:gauss-tools}
The following facts hold for an $n$-dimensional function field in
characteristic zero.
\begin{enumerate}[label=\textup{(\roman*)}]
\item If $w$ is Gauss on a full tuple $t_1,\ldots,t_n$, then it is
quasi-monomial, and
\begin{equation}\label{val:log-frame}
 A_\eta(w)=w\!\left(
 \frac{\eta}{d\log t_1\wedge\cdots\wedge d\log t_n}\right).
\end{equation}
This identity is compatible with finite field extension and pullback
of the form.
\item Suppose $w$ is quasi-monomial and Gauss on a partial tuple $f$.
For nearby assignments of real weights to $f$, there are
quasi-monomial Gauss extensions whose discrepancies and values on any
prescribed finite subset of $L^*$ tend to those of $w$. If $w$ is
centred at $o$, the extensions can be kept centred there.
\item On the finite cone complex of a smooth proper SNC model, the
evaluation of a fixed rational function is continuous, homogeneous,
and rational piecewise linear, including on faces.
\end{enumerate}
\end{lemma}

\begin{proof}
For (i), make a unimodular monomial change in the tuple so that the
first $r_0$ weights are rationally independent and the remaining
weights are zero. Their residues are algebraically independent.
Resolve the supports of the first functions and the rational maps to
$\mathbb P^1$ defined by the others on a proper smooth model. At the
centre, the latter functions are units with independent restrictions,
so the codimension is at most $r_0$. At least $r_0$ boundary components
are needed to express the independent nonzero weights. Thus the centre
is a generic stratum of codimension $r_0$, its boundary parameter
weights are independent, and the valuation is the corresponding
monomial valuation. Indeed different parameter exponents cannot tie,
and the valuation of sufficiently high powers of the maximal ideal
tends to infinity. The vertical exponent matrix of the first
functions has nonzero determinant, and the horizontal residue
differentials of the remaining functions are independent. Their log
wedge is therefore a unit log volume at the stratum, proving
\eqref{val:log-frame}. Applying the same calculation upstairs gives
the finite-extension assertion, consistently with the ramification
formula for form-discrepancies.

For (ii), complete $f$ to a full homogeneous-independent tuple $t$.
This is possible because the graded field of a quasi-monomial
valuation has transcendence degree $n$, by the Abhyankar equality.
Change the prescribed weights of $f$, keep the other tuple weights
fixed, and use Gauss valuations on $k_0(t)$. The field $L$ is finite
over this rational field. For an algebraic element, its possible
extension values are the negatives of Newton-polygon slopes of its
minimal polynomial, with multiplicities. These form a finite multiset
depending continuously on the coefficient values. One may see this
without completeness by extending to an algebraic closure, factoring
the polynomial, and using multiplicativity of Newton polygons.

Several specified values must be tracked by one extension. For a
finite list of elements, choose integer powers so that the value of
their product separates all distinct numerical tuples in the finite
Cartesian product of their candidate value sets at the starting
point. Choose a root of the product's minimal polynomial with value
tending to the prescribed product value. The resulting embedding of
the product field into the algebraic closure extends to $L$. Under
this single embedding every individual value belongs to its candidate
set. Separation forces all those values to converge to the prescribed
ones. Include in the list the coefficient of $\eta$ in the full log
frame, and the functions $h_i$. Formula~\eqref{val:log-frame} gives
discrepancy convergence, and positivity of the $h_i$ preserves the
centre. This argument asserts no uniform size of the neighbourhood.

For (iii), at a generic stratum write numerator and denominator in
the completed regular local ring with a coefficient field and the
boundary parameters. The minimum of the occurring exponent weights
is determined by finitely many exponents, by the Noetherian property
of monomial ideals. This proves rational piecewise linearity. Setting
some weights to zero gives the same evaluation at the containing
stratum. More explicitly, monomial cutoff ideals in the remaining
parameters are primary for the prime generated by those parameters,
and localizing to that stratum and contracting leaves them unchanged.
The same assertion is visible in the power-series completion. This
proves continuity along every face.
\end{proof}

The following form-boundary version uses the transverse-divisor and
toroidal-retraction argument of Jonsson--Musta\c t\u a
\cite[Proposition~5.1, Lemma~5.3 and Corollary~5.4]{JonssonMustata}.
We include the argument for the precise strictness needed here.

\begin{lemma}\label{val:strict-retraction}
Let $(Q,D)$ be a smooth proper model with reduced SNC boundary
containing $\Supp\operatorname{div}_Q(\eta)$. Retraction of a
quasi-monomial valuation $w$ to the boundary cone complex decreases
$A_\eta$ by $A_{Q,D}(w)\ge0$, and the decrease is strict unless $w$
already belongs to that complex.
\end{lemma}

\begin{proof}
The discrepancy formula on a log-smooth model separates into the
boundary-linear contribution and $A_{Q,D}(w)$. Retraction retains
the former. For the equality statement, subdivide the orthant fan
near the centre by rational hyperplanes respecting the smallest
rational subspace containing the boundary-weight vector, and then
take a regular toric refinement. Such subdivisions are realized
locally by toroidal modifications: boundary equations extend to
smooth, or etale, coordinates, so the construction pulls back from
coordinate-axis charts. The reduced pair remains crepant.

The lifted centre lies on the open toric stratum specified by the
relative-interior cone of the weight vector. Weight-zero chart
monomials are units there. The positive normal weights are
rationally independent, since their cone lies in the smallest
rational subspace just chosen. If the centre is not the generic
point of that stratum, choose, near its generic point, a smooth
divisor through the centre transverse to the boundary. Adding it
to the boundary subtracts a strictly positive value from the
discrepancy, while the resulting smooth SNC pair is lc. Thus
$A_{Q,D}(w)>0$. If the centre is generic, independent positive
parameter weights force the monomial valuation, as in
Lemma~\ref{val:gauss-tools}.

The monomial valuations on the refined charts are precisely those
on the subdivided original complex. Indeed the new normal parameters
are toric monomials, and the weight-zero ratios of old boundary
monomials restrict to characters of the torus stratum, with only
the monomial lattice relations at its generic point. The monomial
calculation therefore pulls back unchanged from the orthant charts.
This identifies the equality case on the original complex.
\end{proof}

\subsection{A selected optimal vertex}

Fix centred regular eigenfunctions $f=(f_1,\ldots,f_r)$ and positive
weights $b=(b_1,\ldots,b_r)$ which are linearly independent over $\Q$.
Suppose a centred quasi-monomial extension assigning these weights
exists. Independent weights make every such extension Gauss on $f$.
We minimize
\begin{equation}\label{val:ratio}
 \Phi(w)=\frac{A_\eta(w)}{s(w)}
\end{equation}
over these extensions.

\begin{proposition}\label{val:optimization}
The minimum is attained on a finite rational polyhedral subdivision
of the SNC cone complex of a fixed smooth model. There are finitely
many invariant regular functions, called detectors, such that
maximizing the sum of their values among the minimizers permits the
choice of a vertex. At this selected vertex,
\begin{equation}\label{val:coefficient-vectors}
 w(h)=m_h\cdot b\quad(h\in L^*),\qquad
 A_\eta(w)=e\cdot b,\qquad s(w)=g_s\cdot b,
\end{equation}
where all coefficient vectors are rational and have bounded
denominators for this fixed vertex. Its rational rank is $r$.
For an active generator, meaning $w(h_i)=s(w)v_i$, one has
$m_{h_i}=v_i g_s$.

For fixed germ and tuple, selection is described by finitely many
piecewise semialgebraic branches as $b$ varies. On a segment starting
at independent weights, these branches have a finite partition into
intervals on which the selected ratio is rational-smooth. At
dependent weights the polyhedral prescription is initially relaxed,
without an extra Gauss condition. A compact limit of bounded-ratio
selected vertices from independent positive weights remains a
centred quasi-monomial Gauss extension at the limiting weights.
\end{proposition}

\begin{proof}
Resolve the supports of $h_i,f_j,\eta$ on a smooth proper model
$(Q,D)$. Retraction preserves the values of these functions and $s$.
It therefore remains in the prescribed fibre, and preserves the
Gauss condition because $b$ is independent. By
Lemma~\ref{val:strict-retraction}, an attained minimizer must lie
on the complex.

To construct detectors, take local rational parameter equations for
the boundary components at the generic stratum of each cone. Each
parameter is algebraic over $L^G$. Take a polynomial equation for
it and include invariant regular numerators and denominators for
all its nonzero coefficients. Such presentations exist: an
invariant rational function acquires an invariant regular denominator
by multiplying a denominator by its other group translates. Once
the detector values are fixed, Newton root values allow only
finitely many values for each parameter. Since there are finitely
many cones, simultaneous detector fibres on the complex are finite.

Subdivide so that the evaluations in question and $s$ are linear.
On the closed cones cut out by $w(h_i)\ge0$, every nonzero point
has $A_\eta>0$. To check this, first use rational rays, which are
divisorial and have affine centres because the $h_i$ include affine
generators; klt gives positivity. Linearity on the finitely many
rational cones then gives positivity and properness on each cone.
For $s>0$, \eqref{val:comparison} gives
\begin{equation}\label{val:s-upper}
 s\le\min_j\frac{b_j}{v(f_j)}.
\end{equation}
Thus a bounded sublevel of $\Phi$ has bounded $A_\eta$. A limit
cannot have $s=0$: boundedness of $\Phi$ would force $A_\eta=0$,
hence the zero valuation, contrary to $w(f)=b\ne0$.
The sublevel is therefore compact, and both the minimum and the
secondary maximum are attained.

On each cut polyhedron the minimum locus is the compact face
$A_\eta-\Phi_{\min}s=0$, disjoint from $s=0$. Maximizing the linear
detector sum on that face allows a vertex of the original cut
polyhedron. Its coordinates solve rational linear equations with
right side linear in $b$. Hence they are rational-linear in $b$.
All function values are integral combinations of finitely many
parameter weights at its stratum; this proves the common
denominator assertion. Since the values of the $f_j$ are the
independent $b_j$, the rational rank is exactly $r$. The active
generator identity follows by independence of $b$.

There are only finitely many vertex formulas and polyhedral
feasibility conditions, including $s>0$. Comparisons of their
ratios and then detector sums are semialgebraic. Restriction to a
segment therefore has a finite interval partition. On a valid
branch the denominator is positive, so its ratio is rational-smooth.
For the last assertion, take a convergent subsequence in the fixed
finite complex. The compactness argument keeps $s>0$. For every
fixed polynomial $P$ in the tuple, each independent-weight valuation
satisfies
\[
 w(P(f))=\min_{\alpha:c_\alpha\ne0}\alpha\cdot b,
 \qquad P=\sum_\alpha c_\alpha X^\alpha.
\]
Both sides are continuous on the complex by
Lemma~\ref{val:gauss-tools}. The single limit valuation consequently
satisfies this equality for every polynomial, even when monomial
weights tie. No further subsequence depending on $P$ is needed.
It is therefore Gauss, and admits the approximate continuation of
Lemma~\ref{val:gauss-tools}.
\end{proof}
